% Source problem: public/notes/物工/問題/2018ap_exam.pdf, pages 11–12.
\documentclass[11pt,a4paper]{article}
\usepackage[margin=24mm]{geometry}
\usepackage{fontspec}
\usepackage{xeCJK}
\setCJKmainfont{HaranoAjiMincho-Regular.otf}
\usepackage{amsmath,amssymb,tikz}
\setlength{\parindent}{1em}
\setlength{\parskip}{0.35em}
\title{物理工学専攻 院試解答集\\2018年 物理学II 第3問}
\author{政岡凜太郎}
\date{}
\begin{document}
\maketitle

\section*{[1.1] 入射波の磁場}
非磁性媒質では$\nabla\times\boldsymbol E=-\mu_0\partial_t\boldsymbol H$である。入射電場は$y$方向、波数は$-z$方向で、$E_y^{(i)}=E_0^{(i)}e^{-i(k_0z+\omega t)}$、$k_0=\omega n_0/c$である。したがって
\begin{equation}
 -\frac{\partial E_y^{(i)}}{\partial z}=-\mu_0\frac{\partial H_x^{(i)}}{\partial t}
 \quad\Longrightarrow\quad
 \boxed{H_0^{(i)}=+\frac{n_0}{c\mu_0}E_0^{(i)}}.
\end{equation}
ここで$-\partial_z E_y$は$\nabla\times\boldsymbol E$の$x$成分である。同符号の$E_y,H_x$は$\boldsymbol E\times\boldsymbol H$が$-z$方向を向くこととも整合する。

\section*{[1.2--1.3] 界面条件と振幅反射係数}
表面電流がないので、境界面$z=0$で接線成分が連続する。時刻$t=0$の各振幅を用いれば
\begin{equation}
 E_0^{(i)}+E_0^{(r)}=E_0^{(t)},\qquad
 H_0^{(i)}+H_0^{(r)}=H_0^{(t)}.
\end{equation}
反射波は$+z$方向へ進むので$H_0^{(r)}=-n_0E_0^{(r)}/(c\mu_0)$、透過波は$-z$方向へ進むので$H_0^{(t)}=+n_gE_0^{(t)}/(c\mu_0)$である。従って
\begin{equation}
 n_0(E_0^{(i)}-E_0^{(r)})=n_gE_0^{(t)}.
\end{equation}
これと電場条件を解いて
\begin{equation}
 \boxed{r_0=\frac{E_0^{(r)}}{E_0^{(i)}}=\frac{n_0-n_g}{n_0+n_g}}.
\end{equation}

\section*{[2.1] 各層の電場と磁場}
位相因子$e^{-i\omega t}$のもとで、$E_l^{(-)}$の波は$-z$方向、$E_l^{(+)}$の波は$+z$方向へ進む。Maxwell方程式から符号をそれぞれ定めると
\begin{equation}
 \boxed{H_l^{(-)}=+\frac{n_l}{c\mu_0}E_l^{(-)},\qquad
 H_l^{(+)}=-\frac{n_l}{c\mu_0}E_l^{(+)}}.
\end{equation}

\section*{[2.2] 境界での漸化式}
境界面$l-1$で電場と磁場を連続させる。層$l$側では距離$d_l$だけ伝搬した位相が付くので
\begin{align}
 E_{l-1}^{(-)}+E_{l-1}^{(+)}
 &=E_l^{(-)}e^{-i\Delta_l}+E_l^{(+)}e^{i\Delta_l},\\
 n_{l-1}(E_{l-1}^{(-)}-E_{l-1}^{(+)})
 &=n_l(E_l^{(-)}e^{-i\Delta_l}-E_l^{(+)}e^{i\Delta_l}),
\end{align}
ただし$\Delta_l=n_l\omega d_l/c$である。2式を割り合わせれば問題文の漸化式となり
\begin{equation}
 \boxed{\alpha_l=\frac{n_l}{n_{l-1}}}.
\end{equation}

\section*{[2.3] 四分の一波長多層膜}
空気からガラスへ向かう層の順序と終端を図示する。
\begin{center}
\begin{tikzpicture}[x=1cm,y=0.48cm]
 \draw[thick] (0,0) rectangle (5,0.8);
 \node at (2.5,0.4) {空気 $n_0$};
 \draw[thick] (0,0) rectangle (5,-0.8);
 \node at (2.5,-0.4) {$L$ ($n_L$)};
 \draw[thick] (0,-0.8) rectangle (5,-1.6);
 \node at (2.5,-1.2) {$H$ ($n_H$)};
 \node at (2.5,-2.9) {$\vdots$};
 \draw[thick] (0,-3.5) rectangle (5,-4.3);
 \node at (2.5,-3.9) {$L$ ($n_L$)};
 \draw[thick] (0,-4.3) rectangle (5,-5.1);
 \node at (2.5,-4.7) {$H$ ($n_H$)};
 \draw[very thick] (0,-5.1) rectangle (5,-5.9);
 \node at (2.5,-5.5) {ガラス $n_g$};
 \draw[->,thick] (5.8,0.5) -- (5.8,-0.5) node[midway,right] {入射};
 \draw[->,thick] (6.9,-0.5) -- (6.9,0.5) node[midway,right] {反射};
\end{tikzpicture}
\par\small 図は媒質の順序を示す模式図であり、層の高さは物理的な厚さの比を表さない。
\end{center}
境界での場の比を$Y=c\mu_0H_x/E_y$とする。層$l$の下端$z_l$では、共通時間因子を省いて$E_y=E_l^{(-)}+E_l^{(+)}$、$c\mu_0H_x=n_l(E_l^{(-)}-E_l^{(+)})$なので
\begin{equation}
 Y(z_l)=n_l\frac{E_l^{(-)}-E_l^{(+)}}{E_l^{(-)}+E_l^{(+)}}.
\end{equation}
上端$z_{l-1}=z_l+d_l$では、位相厚さ$\Delta_l=\pi/2$により二つの振幅の位相因子がそれぞれ$-i,+i$となる。従って
\begin{align}
 E_y(z_{l-1})&=-i(E_l^{(-)}-E_l^{(+)}),\\
 c\mu_0H_x(z_{l-1})&=-in_l(E_l^{(-)}+E_l^{(+)}),\\
 Y(z_{l-1})&=n_l\frac{E_l^{(-)}+E_l^{(+)}}{E_l^{(-)}-E_l^{(+)}}
 =\frac{n_l^2}{Y(z_l)}.
\end{align}
ガラス内は透過波だけなので終端の比は$n_g$である。ガラスから空気へ遡ると一組の層は$H$、$L$の順になり、場の比は$Y\mapsto n_H^2/Y\mapsto(n_L^2/n_H^2)Y$と変換される。これを$N$組繰り返すと
\begin{equation}
 Y_{\mathrm{in}}=n_g\left(\frac{n_L}{n_H}\right)^{2N}.
\end{equation}
空気側の入射・反射振幅では$c\mu_0H_x/E_y=n_0(1-r_1)/(1+r_1)$だから
\begin{equation}
 \boxed{r_1=\frac{n_0-Y_{\mathrm{in}}}{n_0+Y_{\mathrm{in}}}
 =\frac{n_0n_H^{2N}-n_gn_L^{2N}}{n_0n_H^{2N}+n_gn_L^{2N}}}.
\end{equation}
設計波長で反射を消す条件は$r_1=0$、すなわち$Y_{\mathrm{in}}=n_0$であり、
\begin{equation}
 \boxed{\frac{n_L}{n_H}=\left(\frac{n_0}{n_g}\right)^{1/(2N)}}.
\end{equation}
このとき積層の入口での場の比が空気の進行波の比に一致する。
\end{document}
